% =====================================================================
%  Chapter 18 review -- the 14-station revision contract.
%  Revision-only: tables, one-liners, pointers back via \secref/\thmref/
%  \propref/\defref/\coderef/\figref/\probref.  No re-derivations.
%  Summit: CF 1181D Irrigation -- water-fill the min-count host city, find
%  the level with a monotone count-space search and name the (offset)-th
%  active city with the frequency-BIT k-th lift (lst:c14-kth).  Samples
%  verified live against the CF statement:
%    (1) n=6 m=4 q=10  a=3 1 1 1 2 2  k=7..16 -> 4 3 4 2 3 4 1 2 3 4
%    (2) n=4 m=5 q=4   a=4 4 5 1      k=15,9,13,6 -> 5 3 3 3
%  Exports prob:c14-summit, lst:c14-kit, lst:c14-summit.
% =====================================================================

\chaptersection{Chapter 18 Summary}

One sentence: \emph{a Fenwick tree is an array that secretly stores a binary
tree over the index axis, so a single \(O(\log n)\) walk answers a prefix sum
and a second walk folds in a point update -- and every question Chapter~16 and
Chapter~17 kept handing back as "just take a prefix" (a range sum, a \(k\)-th
live item, a subtree profile swept offline) is one \(\mathit{lowbit}\) hop
sequence away.} Chapter~5 flattened a tree into an array; this chapter gives
that array a \emph{live prefix engine}. Each row is a pointer, not a proof.

\begin{center}
\renewcommand{\arraystretch}{1.25}
\begin{tabular}{@{}p{0.22\linewidth} p{0.50\linewidth} p{0.20\linewidth}@{}}
\toprule
\textbf{Result} & \textbf{Statement} & \textbf{Where} \\
\midrule
Responsibility interval & Slot \(i\) owns \([\,i-\mathit{lowbit}(i)+1,\ i\,]\);
\(\mathit{lowbit}(i)=i\mathbin\&(-i)\). & \defref{c14-interval} \\
Prefix walk & \(\sum_{1}^{i}A\) is read by jumping \(i\!\to\!i-\mathit{lowbit}(i)\);
\(\bigO{\log n}\). & \thmref{c14-query}, \coderef{c14-query} \\
Point-update walk & Add \(\delta\) to \(A[i]\) by pushing it into every interval
covering \(i\): \(i\!\to\!i+\mathit{lowbit}(i)\). & \coderef{c14-update} \\
Linear build & Hand each partial block to the slot its block extends: \(\bigO{n}\)
build, no updates. & \propref{c14-linear}, \coderef{c14-build} \\
Difference BIT & Two point updates \(+d@l,\ -d@(r{+}1)\) make a range add; a
prefix read is the point value. & \defref{c14-diff}, \coderef{c14-diff-build} \\
Two-BIT range sum & \(\mathit{prefix}(x)=(x{+}1)Q_1(x)-Q_2(x)\) turns range-add
into range-sum. & \thmref{c14-twobit}, \coderef{c14-twobit} \\
Compression & Any value/coordinate set becomes ranks \(1..n\); the BIT only needs
a dense index axis. & \propref{c14-compress}, \coderef{c14-compress} \\
Prefix-min caveat & A min-BIT is \emph{append-only}: valid while values never
decrease back out. & \propref{c14-prefixmin}, \coderef{c14-min} \\
Order statistic & Binary-lift down the tree for the smallest index whose prefix
count reaches \(k\): the \(k\)-th live item. & \thmref{c14-lift}, \coderef{c14-kth} \\
Inversion sweep & A right-to-left frequency-BIT pass counts "smaller after"
online; Ch.~1 owns the merge-sort twin. & \propref{c14-inversion}, \coderef{c14-inv} \\
Euler \(+\) BIT & Ch.~5's \(\tin/\tout\) flat array driven by a BIT: subtree
add/sum by interval, path reads by enter\(+\)1/leave\(-\)1. & \propref{c14-eulerrange}, \propref{c14-path} \\
Dominance & Sort one axis, count the other with a BIT: offline 2D
\(\#\{p_x\le q_x,\ p_y\le q_y\}\). & \propref{c14-dom}, \coderef{c14-dom} \\
PBS over BIT & Parallel binary search whose predicate is a BIT sweep; a BIT only
appends along the sweep, so \emph{no rollback}. & \propref{c14-pbs}, \coderef{c14-pbs} \\
\bottomrule
\end{tabular}
\end{center}

\medskip
{\small\textbf{Algorithms from this chapter.}}
\begin{center}
\renewcommand{\arraystretch}{1.2}
\begin{tabular}{@{}p{0.27\linewidth} l l >{\raggedright\arraybackslash}p{0.26\linewidth}@{}}
\toprule
\textbf{Algorithm} & \textbf{Prep} & \textbf{Per op} & \textbf{Code} \\
\midrule
Prefix sum / point add & \(\bigO{n}\) build & \(\bigO{\log n}\) & \coderef{c14-query}, \coderef{c14-update} \\
CSES 1648 (Dynamic Range Sum) & \(\bigO{n}\) & \(\log n\) & \coderef{c14-1648-soln} \\
Difference BIT (range add / point read) & \(\bigO{n}\) diff & \(\log n\) & \coderef{c14-diff-build} \\
CSES 1651 (Range Update) & \(\bigO{n}\) & \(\log n\) & \coderef{c14-1651-soln} \\
Two-BIT (range add / range sum) & \(\bigO{n}\) & \(\log n\) & \coderef{c14-twobit} \\
Coordinate compression & \(n\log n\) sort & \(\log n\) & \coderef{c14-compress} \\
Frequency \(k\)-th by lifting & \(\bigO{n}\) & \(\log n\) select & \coderef{c14-kth} \\
Inversion / smaller-after & one sweep & \(\log n\)/elem & \coderef{c14-inv} \\
CF 61E (Enemy is weak) & sort ranks & \(\log n\) & \coderef{c14-61-soln} \\
Euler subtree add/sum & \(\bigO{n}\) DFS & \(\log n\) & \coderef{c14-subtree} \\
CF 383C (Propagating tree) & \(\bigO{n}\) Euler & \(\log n\) & \coderef{c14-383-soln} \\
2D dominance count & \(n\log n\) sort & \(\log n\) & \coderef{c14-dom} \\
Nested ranges (CSES 2169) & sort left ends & \(\log n\) & \coderef{c14-nested} \\
PBS over a BIT & \(\bigO{\log R}\) sweeps & \(\log n\) & \coderef{c14-pbs} \\
CF 1181D (Summit, Irrigation) & count + sort & \(\log n\) & \coderef{c14-summit} \\
\bottomrule
\end{tabular}
\end{center}

\chaptersection{Master Decision Flowchart}

Post-study ladder: walk top to bottom, stop at the first branch that matches.
It decides \emph{which BIT flavour} the problem wants and hands off what this
chapter does \emph{not} own: arbitrary-range min with deletions (Ch.~19 segtree),
the merge-sort inversion twin (Ch.~1), lazy propagation (Ch.~20).

\begin{lstlisting}[language={}, basicstyle=\ttfamily\small, numbers=none, frame=single]
  You need a live PREFIX over an index axis.  Which flavour?
  |
  |-- point update + range sum, values combine by + (invertible)?
  |     +-- plain BIT (18.1, lst:c14-query + lst:c14-update). Build O(n)
  |        with lst:c14-build. Range sum = prefix(r)-prefix(l-1).
  |
  |-- add d to A[l..r], but you only ever READ a single point?
  |     +-- difference BIT (18.2, lst:c14-diff-build): +d@l, -d@(r+1),
  |        point read = prefix. Two updates per range (prop:c14-diff).
  |
  |-- add d to A[l..r] AND sum A[l..r]?
  |     +-- two-BIT identity (18.2, lst:c14-twobit) OR lazy segtree.
  |        If you also need min / non-invertible combine -> Ch.19/20 (c13/c15).
  |
  |-- values are huge / sparse, only their ORDER matters?
  |     +-- coordinate compression first (18.2, lst:c14-compress), then any BIT.
  |
  |-- "k-th live item / select / how many below x"?
  |     +-- frequency BIT + binary lifting (18.3, lst:c14-kth). Prefix count
  |        is monotone in the index, so lift down the tree.
  |
  |-- count inversions / smaller-after / greater-before, ONLINE?
  |     +-- one right-to-left frequency sweep (18.3, lst:c14-inv). Merge-sort
  |        (Ch.1, c01) is the offline divide-and-conquer twin -- same bound.
  |
  |-- min over a growing prefix, never pulled back out?
  |     +-- prefix-min BIT (18.2, lst:c14-min), APPEND-ONLY (prop:c14-prefixmin).
  |        Arbitrary decrement -> use a segtree; the caveat is a correctness bug.
  |
  |-- subtree add / subtree sum, or root-path reads, on a TREE?
  |     +-- Euler-flatten (Ch.5, c24) then drive tin/tout with a BIT
  |        (18.4, lst:c14-subtree; path via enter+1/leave-1, lst:c14-path).
  |
  |-- 2D pairs / nested ranges / "how many points dominate q"? offline OK?
  |     +-- sort ONE axis, BIT the other (18.5, lst:c14-dom, lst:c14-nested).
  |
  |-- answer is a THRESHOLD (least x passing a monotone BIT test), many queries?
  |     +-- parallel binary search over one shared BIT (18.5, lst:c14-pbs);
  |        the BIT only appends along the sweep, so NO rollback (Ch.17 did).
  |
  |-- "k-th value in a STATIC range [l..r]"?
        +-- offline value-sweep / PBS over a frequency BIT (18.3 + 18.5); the
           persistent segtree is Ch.19's answer -- this is the offline route.
\end{lstlisting}

\chaptersection{Common Mistakes Checklist}

\begin{itemize}[nosep, leftmargin=1.4em]
  \item \textbf{\(\mathit{lowbit}\) on a signed/zero index.} \(\mathit{lowbit}(i)=i\&(-i)\)
        is only meaningful for \(i\ge 1\); a 0-based BIT loop spins forever
        (\(0+\mathit{lowbit}(0)=0\)). Fenwick is \emph{1-indexed} (\defref{c14-interval}).
  \item \textbf{Query walk vs update walk swapped.} Prefix sum does
        \(i-=\mathit{lowbit}(i)\); point add does \(i+=\mathit{lowbit}(i)\). One
        sign flips the tree from a reader into a writer (\(\coderef{c14-query}\)
        vs \(\coderef{c14-update}\)).
  \item \textbf{Building by \(n\) updates (\(\bigO{n\log n}\)) when \(\bigO{n}\)
        is on offer.} The linear build hands each block forward
        (\(\coderef{c14-build}\)); on \(n=10^6\) the \(\log n\) build is the TLE.
  \item \textbf{Difference BIT missing the \(r+1\) update.} Range add is
        \emph{two} point writes; omitting \(-d@(r{+}1)\) leaks the delta past the
        right end (\(\defref{c14-diff}\)).
  \item \textbf{Treating a min-BIT as deletable.} Prefix-min is append-only; a
        value that later decreases cannot be un-committed
        (\(\propref{c14-prefixmin}\)). That direction is a segtree.
  \item \textbf{Lifting search off-by-one.} The \(k\)-th test is
        \(\mathit{acc}+t[\mathit{next}] < k\) (strict), then return
        \(\mathit{pos}+1\); a \(\le\) returns the wrong neighbour
        (\(\thmref{c14-lift}\), \(\coderef{c14-kth}\)).
  \item \textbf{Forgetting compression on a sparse axis.} A BIT over values up to
        \(10^9\) is not the problem; you must first map values to ranks
        (\(\coderef{c14-compress}\)).
  \item \textbf{\(\tin/\tout\) convention drift.} Ch.~5's \(\tout\) is inclusive;
        subtree sum is \(\mathit{query}(\tout)-\mathit{query}(\tin-1)\) and range
        add writes at \(\tout+1\). Mixing inclusive/exclusive mis-answers every
        subtree (\(\propref{c14-eulerrange}\)).
  \item \textbf{64-bit sums.} Prefix sums and inversion counts reach
        \(\bigO{n^2}\)/\(\bigO{n\cdot\text{val}}\); the tree array is
        \code{long long} the moment values or counts grow (\(\coderef{c14-query}\)).
  \item \textbf{Re-rolling a BIT in PBS.} A BIT sweep is monotone in the sorted
        axis, so PBS needs no undo stack -- unlike Ch.~17's rollback DSU. Copying
        the rollback habit here is pure wasted constant (\(\propref{c14-pbs}\)).
\end{itemize}

\chaptersection{WA Clinic}

Verdicts are symptoms; the causes are the short list this chapter drilled.
Diagnose top-down -- every leaf names a listing.

\begin{figure}[H]
\centering
\begin{tikzpicture}[diagram,
    v/.style={boxnode, fill=primaryblue!12, draw=primaryblue, text width=22mm, align=center},
    b/.style={boxnode, fill=difforange!12, draw=difforange, text width=12mm},
    l/.style={boxnode, fill=successgreen!10, draw=successgreen, text width=56mm,
              font=\scriptsize, align=left}]
  \node[v] (verdict) {\textbf{Verdict on the judge}};
  \node[b, below left=8mm and 4mm of verdict] (wa) {\textbf{WA}};
  \node[b, below right=8mm and 4mm of verdict] (tle) {\textbf{TLE}};
  \node[l, below right=1mm and 2mm of wa.south] (wa1)
    {0-based loop spins or mis-reads: \(\mathit{lowbit}(0)=0\), or the query walk
     uses \(i+=\mathit{lowbit}\) instead of \(i-=\mathit{lowbit}\)
     (\coderef{c14-query}, \coderef{c14-update})};
  \node[l, below left=1mm and 2mm of wa.south] (wa2)
    {difference BIT missing the \(-d@(r{+}1)\) write; inclusive \(\tout\) used as
     exclusive in a subtree sum (\coderef{c14-diff-build}, \propref{c14-eulerrange})};
  \node[l, below=1mm of tle.south] (tle1)
    {\(\bigO{n\log n}\) build where \(\bigO{n}\) was required; un-compressed axis
     sized \(10^9\) (\coderef{c14-build}, \coderef{c14-compress})};
  \node[l, below left=5mm and 30mm of tle1.north] (tle2)
    {int overflow on prefix sums / inversions -- tree array not
     \code{long long} (\(\coderef{c14-inv}\))};
  \draw[flow] (verdict) -- (wa);
  \draw[flow] (verdict) -- (tle);
  \draw[flow] (wa) -- (wa1);
  \draw[flow] (wa) -- (wa2);
  \draw[flow] (tle) -- (tle1);
  \draw[flow] (tle) -- (tle2);
\end{tikzpicture}
\caption*{WA: the walk went the wrong way, or a range write forgot its partner.
TLE: you built with updates, sized the tree by value, or overflowed to
\code{int}. RE: an un-compressed \(10^9\) index into a fixed array.}
\end{figure}

\chaptersection{Quick Reference Card}

\begin{quickref}[Chapter 18 at a glance]
\textbf{Plain BIT.} \(\mathit{lowbit}(i)=i\&(-i)\).
\code{query(i)}: \(s\!+=\!t[i];\,i-=\!\mathit{lowbit}(i)\) sums \(A[1..i]\).
\code{update(i,d)}: \(t[i]\!+=\!d;\,i+=\!\mathit{lowbit}(i)\) adds to \(A[i]\).
Range sum \(=\) \code{query(r)}\(-\,\)code{query(l-1)}
(\(\thmref{c14-query}\), \(\coderef{c14-query}\)).

\medskip
\textbf{O(n) build.} \(t[i]+=\!A[i]\); let \(j=i+\mathit{lowbit}(i)\); if
\(j\le n\) then \(t[j]+=\!t[i]\) (\(\propref{c14-linear}\), \(\coderef{c14-build}\)).

\medskip
\textbf{Range add / point read.} Difference BIT: write \(+d\) at \(l\), \(-d\) at
\(r{+}1\); the value at \(x\) is \code{query(x)}
(\(\defref{c14-diff}\), \(\coderef{c14-diff-build}\)).

\medskip
\textbf{Range add / range sum.} Two BITs \(Q_1,Q_2\); \(\mathit{prefix}(x)=(x{+}1)Q_1(x)-Q_2(x)\);
a range add is two \(\mathit{internalAdd}\) calls
(\(\thmref{c14-twobit}\), \(\coderef{c14-twobit}\)).

\medskip
\textbf{Order statistic.} Frequency BIT; lift down the tree for the least index
with prefix count \(\ge k\): \(k\)-th live item, strictly-\(<x\) count
(\(\thmref{c14-lift}\), \(\coderef{c14-kth}\)).

\medskip
\textbf{On a tree.} Flatten once (Ch.~5 \(\tin/\tout\)); a subtree is an interval
(BIT range op), a root path is enter\(+\)1/leave\(-\)1
(\(\propref{c14-eulerrange}\), \(\coderef{c14-subtree}\), \(\coderef{c14-path}\)).

\medskip
\textbf{Offline.} Sort one axis, BIT the other for dominance/nested counts
(\(\coderef{c14-dom}\), \(\coderef{c14-nested}\)); wrap that count in a parallel
binary search for a threshold -- no rollback, a BIT only appends
(\(\propref{c14-pbs}\), \(\coderef{c14-pbs}\)).

\medskip
\textbf{Translation.} "sum a range, edit a point" \(\to\) plain BIT. "add to a
range, read a point" \(\to\) difference BIT. "add to a range, sum a range" \(\to\)
two-BIT (or segtree). "which live value is the \(k\)-th? how many below \(x\)?"
\(\to\) frequency lifting. "subtree/path on a tree" \(\to\) Euler \(+\) BIT.
"how many points dominate \(q\)" / "least \(x\) passing a prefix test" \(\to\)
offline sweep, PBS over a BIT.
\end{quickref}

\begin{lstlisting}[caption={Chapter 18 kit: the plain BIT, the range-add/range-sum pair, and the k-th lift, complete and production-shaped}, label={lst:c14-kit}]
struct BIT {                                  // 1-indexed, long long payload
    int n; vector<long long> t;
    BIT(int n):n(n),t(n+1,0){}
    static int lowbit(int i){return i&(-i);}
    void add(int i,long long d){for(;i<=n;i+=lowbit(i))t[i]+=d;}   // point write
    long long pre(int i){long long s=0;for(;i>0;i-=lowbit(i))s+=t[i];return s;} // prefix
    long long range(int l,int r){return pre(r)-pre(l-1);}          // invertible combine
    // O(n) build: hand each partial block forward
    void build(const vector<long long>&a){t.assign(n+1,0);
        for(int i=1;i<=n;i++){t[i]+=a[i];int j=i+lowbit(i);if(j<=n)t[j]+=t[i];}}
    // k-th: least idx whose prefix sum reaches k (monotone counts). hb = 2^floor(log2 n)
    int kth(long long k,int hb){int pos=0;long long acc=0;
        for(int pw=hb;pw>0;pw>>=1){int nx=pos+pw;
            if(nx<=n && acc+t[nx]<k){pos=nx;acc+=t[nx];}} return pos+1;}
};
// range add + range sum with two BITs (thm:c14-twobit)
struct RangeAddSum{int n;BIT q1,q2;RangeAddSum(int n):n(n),q1(n),q2(n){}
    void ia(int i,long long v){q1.add(i,v);q2.add(i,v*(i-1));}
    void range_add(int l,int r,long long v){ia(l,v);ia(r+1,-v);}
    long long pref(int x){return (long long)(x+1)*q1.pre(x)-q2.pre(x);}
    long long range_sum(int l,int r){return pref(r)-pref(l-1);}};
// highest power of two <= n, for kth:   int hb=1;while(hb*2<=n)hb*=2;
\end{lstlisting}

\chaptersection{Write-From-Memory Drills}

Five signature skeletons, each lifted from a listing already in the chapter.
Cover the marked line, reproduce from memory, then check the answer one line
below.

\begin{lstlisting}[caption={Drill 1 -- prefix query + point update (from lst:c14-query / lst:c14-update)}]
int lowbit(int i) { return i & (-i); }
long long query(int i) {                 // sum of A[1..i]
    long long s = 0;
    for (; i > 0; i -= lowbit(i)) s += tree[i];   // WALK DOWN
    return s;
}
void update(int i, long long d) {        // add d to A[i]
    for (; i <= n; i += lowbit(i)) tree[i] += d;  // WALK UP
}
\end{lstlisting}
\noindent{\scriptsize Key: the query walks \emph{down} (\(-\)lowbit) collecting
disjoint responsibility blocks; the update walks \emph{up} (\(+\)lowbit) touching
every block that contains \(i\). Range sum \(=\) \code{query(r)-query(l-1)}.}

\begin{lstlisting}[caption={Drill 2 -- O(n) linear build (from lst:c14-build)}]
for (int i = 1; i <= n; i++) {
    tree[i] += a[i];                        // own slot's value in
    int j = i + lowbit(i);                  // the block that EXTENDS i's
    if (j <= n) tree[j] += tree[i];         // hand the partial block forward
}
\end{lstlisting}
\noindent{\scriptsize Key: each \(t[i]\) already sums \([i-\mathit{lowbit}(i)+1,i]\),
so adding it to \(t[j]\) extends \(j\)'s interval in one pass -- \(\bigO{n}\), not
\(n\) updates (\(\propref{c14-linear}\)).}

\begin{lstlisting}[caption={Drill 3 -- difference BIT + two-BIT range sum (from lst:c14-diff-build / lst:c14-twobit)}]
// range add / POINT read
void range_add_diff(int l, int r, long long d) { add(l, d); add(r + 1, -d); }
long long point_read(int x) { return query(x); }         // prefix == value

// range add / RANGE sum  (two BITs)
void internal_add(int i, long long v) { bit1.add(i, v); bit2.add(i, v * (i - 1)); }
void range_add(int l, int r, long long v) { internal_add(l, v); internal_add(r + 1, -v); }
long long prefix(int x) { return (x + 1) * bit1.query(x) - bit2.query(x); }
long long range_sum(int l, int r) { return prefix(r) - prefix(l - 1); }
\end{lstlisting}
\noindent{\scriptsize Key: forgetting the \(-d@(r{+}1)\) write in the difference
BIT leaks the delta past the right end; the two-BIT form is just the difference
idea weighted by position (\(\thmref{c14-twobit}\)).}

\begin{lstlisting}[caption={Drill 4 -- k-th by binary lifting (from lst:c14-kth)}]
int kth(long long k) {           // least idx with prefix count >= k; hb = 2^floor(log2 n)
    int pos = 0; long long acc = 0;
    for (int pw = hb; pw > 0; pw >>= 1) {
        int nxt = pos + pw;
        if (nxt <= n && acc + tree[nxt] < k) {   // STRICT <, else wrong neighbour
            pos = nxt; acc += tree[nxt];
        }
    }
    return pos + 1;              // first index pushing the prefix to k
}
\end{lstlisting}
\noindent{\scriptsize Key: descent is the \(\bigO{\log n}\) search down the hidden
binary tree; the test is \(\mathit{acc}+t[\mathit{nxt}]<k\) so we stop just short
and return \(\mathit{pos}+1\) (\(\thmref{c14-lift}\)).}

\begin{lstlisting}[caption={Drill 5 -- Euler subtree add / sum (from lst:c14-subtree)}]
// tin/tout from Ch.5 DFS (tout INCLUSIVE); BIT over the preorder flat array
long long subtree_sum(int v)  { return query(tout[v]) - query(tin[v] - 1); }
void subtree_add(int v, long long d) {
    add(tin[v], d);              // difference BIT on the Euler axis
    add(tout[v] + 1, -d);        // close the interval at tout (inclusive!)
}
long long node_value(int v) { return query(tin[v]); }   // point read after adds
\end{lstlisting}
\noindent{\scriptsize Key: a subtree is exactly the interval \([\tin,\tout]\); the
\(+1/-1\) pair is the difference-BIT range add applied on the flattened axis
(\(\propref{c14-eulerrange}\)). Off-by-one lives entirely in the \(\tout\)
convention.}

\chaptersection{Complete Problem Checklist}

Every problem this chapter offers -- 4 fully solved in-section, worked offline
entries, practice, and one Summit. Tick \emph{only} after an unaided solve; a
hinted attempt counts as half.

\begin{center}
\footnotesize
\renewcommand{\arraystretch}{1.25}
\begin{tabular}{@{}c p{0.19\linewidth} l p{0.20\linewidth} p{0.24\linewidth} c@{}}
\toprule
\textbf{\#} & \textbf{Problem} & \textbf{Source} & \textbf{Key technique} & \textbf{What to look for} & \textbf{Done?} \\
\midrule
\multicolumn{6}{l}{\textbf{Section 18.1 -- BIT core}} \\
W18.1 & Dynamic Range Sum Queries & CSES 1648 & \lv{3} plain BIT, set-as-delta & point \emph{set} \(\to\) \(\delta=x-A[i]\); range \(=\) two prefixes & \(\square\) \\
\midrule
\multicolumn{6}{l}{\textbf{Section 18.2 -- range variants}} \\
W18.2 & Range Update Queries & CSES 1651 & \lv{4} difference BIT & two writes per range; point read \(=\) prefix & \(\square\) \\
W18.2b & Range add + range sum & worked & \lv{4} two-BIT identity & \((x{+}1)Q_1-Q_2\); mind the \(r{+}1\) write & \(\square\) \\
\midrule
\multicolumn{6}{l}{\textbf{Section 18.3 -- order statistics}} \\
W18.3 & Enemy is weak & CF 61E & \lv{5} frequency BIT, two sweeps & per-middle product of (greater-left)\(\times\)(smaller-right) & \(\square\) \\
W18.3b & \(k\)-th live item / Hotel & worked & \lv{4} binary lifting & strict-\(<k\) descent; return \(\mathit{pos}+1\) & \(\square\) \\
\midrule
\multicolumn{6}{l}{\textbf{Section 18.4 -- Fenwick on trees}} \\
W18.4 & Propagating tree & CF 383C & \lv{5} Euler + parity difference BIT & depth-parity sign folded into the delta; inclusive \(\tout\) & \(\square\) \\
\midrule
\multicolumn{6}{l}{\textbf{Section 18.5 -- offline sweeps (worked)}} \\
W18.5a & 2D dominance count & worked & \lv{4} sort one axis, BIT the other & compress the second; count \(\le\), not \(<\) & \(\square\) \\
W18.5b & Nested ranges & CSES 2169 & \lv{4} sort left ends & tie-break order decides contain/contained & \(\square\) \\
W18.5c & PBS over a BIT & worked & \lv{5} lockstep search, BIT predicate & no rollback -- the sweep only appends & \(\square\) \\
\midrule
\multicolumn{6}{l}{\textbf{Summit (chapter end)}} \\
W18 & Irrigation & CF 1181D & \lv{5} water-fill level + \(k\)-th active city & count-space search, then lift the offset-th index & \(\square\) \\
\bottomrule
\end{tabular}
\end{center}

\chaptersection{Recommended Study Order}

Five focused days. Day~1 internalises the walks and the linear build; day~2 the
difference and two-BIT range variants plus compression; day~3 the frequency BIT
and lifting; day~4 Fenwick on the Euler array; day~5 the offline sweeps and the
Summit.

\begin{center}
\small
\renewcommand{\arraystretch}{1.25}
\begin{tabular}{@{}llp{0.55\linewidth}@{}}
\toprule
\textbf{Day} & \textbf{Focus} & \textbf{Tasks} \\
\midrule
1 & Read + do 18.1 & re-trace the query/update walks on \(A\) (\figref{c14-array}) by hand; write \coderef{c14-query}/\coderef{c14-update} cold; CSES~1648 end to end \\
2 & Read + do 18.2 & difference cascade on \(A\); CF... none -- re-derive \coderef{c14-twobit}'s \((x{+}1)Q_1-Q_2\); CSES~1651; one compression pass with \coderef{c14-compress} \\
3 & Read + do 18.3 & lift search for the \(k\)-th on the frequency array (\figref{c14-lift}); right-to-left inversion sweep (\coderef{c14-inv}); CF~61E \\
4 & Read + do 18.4 & overlay \(\tin/\tout\) of \(T\) (\figref{c14-euler}) on the BIT; subtree add/sum by hand; CF~383C parity fold \\
5 & Do 18.5 + Summit & dominance count and nested-ranges sweeps (\(\coderef{c14-dom}\), \(\coderef{c14-nested}\)); PBS skeleton \(\coderef{c14-pbs}\); then CF~1181D; walk the flowchart; first 3 mixed \\
\bottomrule
\end{tabular}
\end{center}

\chaptersection{Difficulty Progression}

Tiers are competency gates -- each row names the \emph{skill} the tier unlocks.

\begin{center}
\small
\renewcommand{\arraystretch}{1.3}
\begin{tabular}{@{}lp{0.34\linewidth}p{0.38\linewidth}@{}}
\toprule
\textbf{Tier} & \textbf{Problems} & \textbf{What passing this tier proves} \\
\midrule
\lv{3}--\lv{4} \ \(\bigstar\bigstar\bigstar\bigstar\)
  & CSES~1648, CSES~1651, worked two-BIT
  & you can build a BIT cold, drive the query/update walks, and pick difference vs two-BIT for the right range semantics \\
\lv{4}--\lv{5} \ \(\bigstar\bigstar\bigstar\bigstar\bigstar\)
  & CF~61E, CF~383C, worked dominance/nested
  & you turn an index into a \emph{meaning} (rank, Euler position, sorted axis) and lift the tree for an order statistic \\
\lv{5}{+} \ \(\bigstar\bigstar\bigstar\bigstar\bigstar\)
  & CF~1181D (Summit), CF~1045G, CF~220E, CF~1000F
  & you compose two BIT levels (search a threshold \emph{and} select the item) and harden for \(n,q\le 5\cdot10^5\), \(k\le 10^{18}\) \\
\bottomrule
\end{tabular}
\end{center}

\chaptersection{Reduction Gallery}

Every entry is ``one BIT flavour + one twist.'' Each line collapses a problem
family onto a listing already in the chapter.

\begin{itemize}[nosep, leftmargin=1.4em]
  \item ``sum a range after point edits'' (CSES~1648) \(\to\) plain BIT, set as
        \(\delta\), difference of two prefixes (\(\coderef{c14-query}\), \(\probref{c14-rangesum}\)).
  \item ``add to a range, read a point'' (CSES~1651) \(\to\) difference BIT, two
        writes per range (\(\defref{c14-diff}\), \(\probref{c14-rangeupd}\)).
  \item ``add to a range, sum a range'' (worked) \(\to\) two-BIT identity
        \((x{+}1)Q_1-Q_2\) (\(\thmref{c14-twobit}\), \(\coderef{c14-twobit}\)).
  \item ``\(k\)-th live item / how many \(<x\)'' (Hotel, worked) \(\to\) frequency
        BIT + binary lifting (\(\thmref{c14-lift}\), \(\coderef{c14-kth}\)).
  \item ``decreasing triplet / smaller-after count'' (CF~61E, LC~315) \(\to\) one
        or two frequency sweeps over ranks (\(\propref{c14-inversion}\), \(\probref{c14-triplets}\)).
  \item ``subtree add/sum, root path'' (CF~383C, CSES~1137) \(\to\) Ch.~5 flatten,
        then difference BIT on the Euler axis (\(\propref{c14-eulerrange}\), \(\probref{c14-prop}\)).
  \item ``how many points dominate \(q\) / nested ranges'' (worked) \(\to\) sort one
        axis, BIT the other (\(\propref{c14-dom}\), \(\coderef{c14-dom}\)).
  \item ``least threshold passing a prefix test, many queries'' (CF~1181D) \(\to\)
        PBS over one shared BIT sweep (\(\propref{c14-pbs}\), \(\probref{c14-summit}\)).
\end{itemize}

\chaptersection{Mixed Problem Set}

Ten entries from CF~1800 to CF~2400, all verified against the live problem set
and absent from the rest of the book. None says ``use a Fenwick tree'' --
recognise the flavour from what the problem needs summed, selected, or swept.

\begin{mixedset}
  \practice{Pashmak and Parmida's problem}{CF 459D}
    {https://codeforces.com/problemset/problem/459/D}{CF 1800 L3}
    {Build \(L[i]\)=\#occurrences of \(a_i\) in \(a[1..i]\) and \(R[i]\)=\# in
     \(a[i..n]\) with two map passes, then count pairs \(i<j\) with \(L[i]>R[j]\):
     that is a strictly-greater-before sweep -- insert \(L\) values into a
     frequency BIT keyed by rank and query the complement. The very
     right-to-left count of \(\coderef{c14-inv}\) on the \(R\) array
     (\(\secref{c14-order}\)).}
  \practice{Closest Equals}{CF 522D}
    {https://codeforces.com/problemset/problem/522/D}{CF 2000 L4}
    {Over queries \([l,r]\), the smallest \(|x-y|\) with \(a_x=a_y\), else \(-1\).
     Unique candidate keys are the \((p,\text{next}[p])\) equal-value pairs; sweep
     \(l\) downward and commit each pair's gap at coordinate \(\text{next}[p]\) into
     a \emph{prefix-min} BIT (append-only \(\Rightarrow\) legal), answer
     \(\min\) over \([l,r]\) by a range read. \(\coderef{c14-min}\) with the
     append caveat from \(\propref{c14-prefixmin}\) (\(\secref{c14-range}\)).}
  \practice{Yaroslav and Divisors}{CF 301D}
    {https://codeforces.com/problemset/problem/301/D}{CF 2200 L5}
    {For each query \([l,r]\), count pairs \((i,j)\), \(l\le i\le j\le r\), with
     \(a_i\mid a_j\) (or vice versa). Pre-list all valid \((i,j)\) pairs, sort them
     and the queries by right end, sweep with a BIT over left ends -- the exact
     nested-ranges count of \(\coderef{c14-nested}\) applied to pair intervals
     (\(\secref{c14-offline}\)).}
  \practice{Radio stations}{CF 762E}
    {https://codeforces.com/problemset/problem/762/E}{CF 2200 L5}
    {Count pairs of stations with \(\min(r_i,r_j)\ge|x_i-x_j|\) and
     \(|f_i-f_j|\le k\). Sort stations by range endpoint descending and activate
     them into a BIT over the (compressed) frequency axis; for a new station the
     admissible partners are a \(2k{+}1\) window of frequencies, read as frequency
     \(\times\) coordinate-prefix counts. The dominance sweep of
     \(\coderef{c14-dom}\) with a \(\coderef{c14-freq}\) inside
     (\(\secref{c14-offline}\)).}
  \practice{The Number Games}{CF 980E}
    {https://codeforces.com/problemset/problem/980/E}{CF 2200 L5}
    {Keep \(n-k\) tree nodes, connected and containing root \(1\), maximising
     \(\sum 2^i\) over kept labels -- value dominance forces a greedy over
     labels \(v=n..1\). Including \(v\) commits its whole path to the already-kept
     structure; a BIT over the Euler order counts how many still-claimable path
     nodes a subtree holds, bounding the pick. Euler-array driving from
     \(\coderef{c14-subtree}\) / \(\coderef{c14-path}\) (\(\secref{c14-trees}\)).}
  \practice{AI robots}{CF 1045G}
    {https://codeforces.com/problemset/problem/1045/G}{CF 2200 L5}
    {Count pairs of robots that see each other (\(|\Delta\text{position}|\) within
     both sight radii) with \(|q_i-q_j|\le K\). Group by IQ and two-pointer over
     the \(K\)-window; within a merge of two IQ groups, sort by position and count
     mutual-visibility pairs with one BIT over the compressed \(x\pm r\) endpoints
     -- a straight 2D dominance \(\coderef{c14-dom}\) after
     \(\coderef{c14-compress}\) (\(\secref{c14-offline}\)).}
  \practice{Fenwick Tree}{CF 1967C}
    {https://codeforces.com/problemset/problem/1967/C}{CF 2300 L5}
    {A game read directly off the \(\mathit{lowbit}\) structure: who wins depends on
     the parity of the number of set bits reached by alternating
     \(i\!\to\!i-\mathit{lowbit}(i)\) descents. Internalise \(\coderef{c14-query}\)'s
     walk and \(\defref{c14-interval}\)'s responsibility intervals until the
     \(\bigO{\log n}\) hop count is reflex (\(\secref{c14-core}\)).}
  \practice{Little Elephant and Inversions}{CF 220E}
    {https://codeforces.com/problemset/problem/220/E}{CF 2400 L5}
    {Count subarrays whose \emph{removal} leaves \(\le k\) inversions. Two pointers
     on the removed window while maintaining the surviving inversion count -- add
     and remove an element's inversion contribution through a frequency BIT over
     ranks. The online count of \(\coderef{c14-inv}\), made sliding, with
     \(\coderef{c14-freq}\) as the live structure (\(\secref{c14-order}\)).}
  \practice{Little Elephant and Tree}{CF 258E}
    {https://codeforces.com/problemset/problem/258/E}{CF 2400 L5}
    {Subtree colour updates and ``how many distinct colours fully cover subtree
     \(v\)'' queries. Flatten with \(\tin/\tout\) (Ch.~5), then a range-add /
     point-read difference BIT per colour tracks coverage depth, and the offline
     query ordering is the sweep of \(\coderef{c14-dom}\) over the Euler axis
     (\(\secref{c14-trees}\), \(\coderef{c14-diff-build}\)).}
  \practice{One Occurrence}{CF 1000F}
    {https://codeforces.com/problemset/problem/1000/F}{CF 2400 L5}
    {Answer queries \([l,r]\): report any value appearing exactly once, else \(0\).
     Sweep right end \(r\) left-to-right, keep each value's last and second-last
     positions in a BIT over positions (contribute \(+1\) at the newest, remove the
     stale ones); read a point in \([l,r]\) for a singleton. Offline nested-range
     bookkeeping as in \(\coderef{c14-nested}\) (\(\secref{c14-offline}\)).}
\end{mixedset}

\chaptersection{The Summit -- CF 1181D Irrigation}

\begin{notebox}[Why this is the chapter top rope]
Sections 18.3 and 18.5 each own one half: 18.3 turns a prefix \emph{count} into a
\(k\)-th \emph{item} by lifting the tree; 18.5 sorts queries and sweeps one axis
with a BIT. The Summit is both at once, twice: it must first \emph{find a water
level} by searching a monotone count-space prefix function, then \emph{select the}
\(\mathit{offset}\)-th still-short city -- a \(k\)-th lift over a live frequency
BIT -- and it does so for \(q\le 5\cdot10^5\) queries with year numbers up to
\(10^{18}\), so nothing can be simulated year by year. Compress the counts, sort
queries offline, and let the two BIT walks answer each in \(\bigO{\log n}\).
\end{notebox}

\solvedproblem{c14-summit}{Irrigation}{Codeforces 1181D}
  {https://codeforces.com/problemset/problem/1181/D}{CF 2200 L5}
  {water-fill over host counts + offline sorted queries + frequency-BIT k-th lifting over city index}

\begin{problemstatement}[CF 1181D Irrigation]
There are \(m\) cities, numbered \(1..m\). The first \(n\) olympiads are given:
year \(i\) was hosted by city \(a_i\). From year \(n+1\) on, the host is chosen by
a rule: \emph{the city that has hosted the fewest times so far; ties go to the
smallest index}. Each of \(q\) queries gives a year \(k_i\) with
\(n+1\le k_i\le 10^{18}\); report the host city that year.

\textbf{Input.} \(n,m,q\) (\(1\le n,m,q\le 5\cdot10^5\)); then \(n\) integers
\(a_1..a_n\) (\(1\le a_i\le m\)); then \(q\) lines, each \(k_i\)
(\(n+1\le k_i\le 10^{18}\)).

\textbf{Output.} \(q\) integers -- the host city in each queried year.

\textbf{Sample 1.} \(n=6,\ m=4,\ q=10\), \(a=(3,1,1,1,2,2)\),
\(k=7,8,\dots,16\) \(\to\) \texttt{4 3 4 2 3 4 1 2 3 4}.\\
\textbf{Sample 2.} \(n=4,\ m=5,\ q=4\), \(a=(4,4,5,1)\),
\(k=15,9,13,6\) \(\to\) \texttt{5 3 3 3}.
\end{problemstatement}

\solstep{Key Observations}
\begin{enumerate}[nosep,leftmargin=1.6em]
  \item Let \(c[x]\) be how many of the first \(n\) years city \(x\) hosted.
        The rule always lifts a current minimum, so it \emph{equalises} the
        counts: define \(S(L)=\sum_x \max(0,\,L-c[x])\), the number of post-rule
        selections needed to bring \emph{every} city to at least \(L\). \(S\) is
        non-decreasing in \(L\) -- a monotone prefix function.
  \item Write a query as \(j=k-n\), the index of the selection after the rule
        starts (1-based). The chosen city sits in the unique \emph{band}
        \(S(L-1)<j\le S(L)\): precisely the cities with \(c[x]<L\), all standing
        at \(L-1\), picked in increasing index order. So the answer is the
        \(\big(\,j-S(L-1)\,\big)\)-th smallest index among cities with \(c[x]<L\).
  \item Finding \(L\) is a monotone search over a count axis -- \(18.5\)'s PBS /
        binary-search trick, evaluating \(S\) with prefix sums of a
        \emph{frequency-of-count} BIT. Picking the offset-th active city is
        \(18.3\)'s \(k\)-th lift (\(\thmref{c14-lift}\), \(\coderef{c14-kth}\)).
\end{enumerate}

\solstep{Reduction to a Lift over a Sweep}
Two BITs, two axes. \(\mathcal{B}_{\mathit{cnt}}\): a frequency BIT over count
values (how many cities host exactly \(h\) times) with running sums of
\(\#\text{cities}\) and \(\sum\text{counts}\), so \(S(L)=L\cdot C(L)-P(L)\) is a
prefix read. \(\mathcal{B}_{\mathit{idx}}\): a frequency BIT over city indices,
city \(x\) active iff \(c[x]<L\); its \((\,j-S(L-1)\,)\)-th active slot is found
by \(\coderef{c14-kth}\)'s descent. Process queries sorted by \(k\): \(L\) is
non-decreasing across them, so \(\mathcal{B}_{\mathit{idx}}\) only ever
\emph{appends} activations -- exactly the append-only regime that lets
\(\propref{c14-pbs}\)'s PBS skip a rollback. Each query: advance \(L\) while
\(S(L)<j\), activate newly-qualifying cities, then lift the
\((j-S(L-1))\)-th active index. \(\bigO{(n+q)\log n}\) after a \(\bigO{n}\) count.

\solstep{Trace on Sample 1}
\begin{tracebox}[Water-fill on \(n=6,m=4\), \(a=(3,1,1,1,2,2)\)]
Counts \(c=(3,2,1,0)\) for cities \(1..4\). \(S(L)=\sum_x\max(0,L-c[x])\):
\[
{\small
\begin{array}{c|ccccc}
L & 0 & 1 & 2 & 3 & 4\\ \hline
S(L) & 0 & 1 & 3 & 6 & 10
\end{array}}
\qquad
{\small
\begin{array}{c|c|c|c|c|c}
k & j{=}k{-}6 & \text{band }L & \text{batch }(\text{cities }c{<}L) & \text{offset} & \text{host}\\ \hline
7&1&1&\{4\}&1&4\\
8&2&2&\{3,4\}&1&3\\
9&3&2&\{3,4\}&2&4\\
10&4&3&\{2,3,4\}&1&2\\
11&5&3&\{2,3,4\}&2&3\\
12&6&3&\{2,3,4\}&3&4\\
13&7&4&\{1,2,3,4\}&1&1\\
14&8&4&\{1,2,3,4\}&2&2\\
15&9&4&\{1,2,3,4\}&3&3\\
16&10&4&\{1,2,3,4\}&4&4
\end{array}}
\]
Read the batch column as \(\mathcal{B}_{\mathit{idx}}\)'s active set (indices
ascending) and the offset as the lift target; the answers
\(4,3,4,2,3,4,1,2,3,4\) match the official output. The jump from \(j=6\) (offset
\(3\) at level \(3\)) to \(j=7\) (level \(4\), offset resets to \(1\)) is where
city~1 first joins because \(S(3)=6<7\le S(4)\).
\end{tracebox}

\solstep{Implementation}
\begin{lstlisting}[caption={CF 1181D: sort counts, search the water level, lift the offset-th active city}, label={lst:c14-summit}]
#include <bits/stdc++.h>
using namespace std;
int bit_[500005];                        // frequency BIT over city indices
int M, HB;                               // HB = 2^floor(log2 M)
void add(int i,int d){for(;i<=M;i+=i&-i)bit_[i]+=d;}
int kth(long long k){                    // least idx whose prefix >= k
    int pos=0; long long acc=0;
    for(int pw=HB;pw>0;pw>>=1){int nx=pos+pw;
        if(nx<=M && acc+bit_[nx]<k){pos=nx;acc+=bit_[nx];}}
    return pos+1;}
int main(){
    int n,q; scanf("%d%d%d",&n,&M,&q);
    vector<int> cnt(M+1,0);
    for(int i=0,a;i<n;i++){scanf("%d",&a);cnt[a]++;}
    vector<pair<int,int>> ord;           // (count, city), sorted count asc
    for(int x=1;x<=M;x++) ord.push_back({cnt[x],x});
    sort(ord.begin(),ord.end());
    vector<int> cs(M); vector<long long> ps(M+1,0);
    for(int t=0;t<M;t++) cs[t]=ord[t].first;
    for(int t=0;t<M;t++) ps[t+1]=ps[t]+cs[t];      // prefix sum of counts
    auto S=[&](long long L)->__int128{             // selections to reach level L
        long long k=lower_bound(cs.begin(),cs.end(),L)-cs.begin(); // #{cnt<L}
        return (__int128)L*k-ps[k];};
    HB=1; while(HB*2<=M)HB*=2;
    vector<pair<long long,int>> qs(q);
    for(int i=0;i<q;i++){long long k;scanf("%lld",&k);qs[i]={k-n,i};}
    sort(qs.begin(),qs.end());            // by j ascending => L monotone
    vector<int> ans(q); int act=0;        // act: cities with cnt < L activated
    for(auto&qu:qs){
        long long j=qu.first;             // 1-based selection index
        long long lo=1,hi=(long long)2e9;
        while(lo<hi){long long mid=(lo+hi)/2;if(S(mid)>=j)hi=mid;else lo=mid+1;}
        long long L=lo;                   // smallest level with S(L) >= j
        while(act<M && ord[act].first<L){ add(ord[act].second,1); act++; } // append-only
        long long used=(long long)(S(L-1)); // = S(L-1) fits in int64 (<=j)
        long long off=j-used;             // 1-based inside the band
        ans[qu.second]=kth(off);          // off-th smallest active city index
    }
    for(int i=0;i<q;i++) printf("%d ",ans[i]);
}
// Two BIT axes: S(L) is a count-space prefix read (sorted counts + prefix sums);
// kth(off) is 18.3's index-space lift.  Sorting queries by j makes L non-decreasing,
// so activations are appended once -- no rollback, exactly prop:c14-pbs's argument.
// __int128 guards L*k on the level search; j,k <= 1e18 need long long.
\end{lstlisting}

\begin{hardeningbox}[AC \(\longrightarrow\) red: hardening the Summit]
\begin{itemize}[nosep,leftmargin=1.4em]
  \item \textbf{Never simulate years.} \(k_i\le 10^{18}\); the only tractable route
        is the closed-form band \(S(L-1)<j\le S(L)\) with \(j=k-n\). A per-year
        loop is the WA/TLE reflex to kill first.
  \item \textbf{Two axes, one idea.} The level search lives in \emph{count} space
        (\(\mathcal{B}_{\mathit{cnt}}\)); the city pick lives in \emph{index} space
        (\(\mathcal{B}_{\mathit{idx}}\), \(\coderef{c14-kth}\)). Conflating them is
        the classic off-by-one that shifts every answer by one city.
  \item \textbf{Activation is append-only.} Because \(L\) never decreases across
        sorted queries, cities enter \(\mathcal{B}_{\mathit{idx}}\) once, in count
        order. Do \emph{not} rebuild the BIT per query -- hoist it, exactly the
        no-rollback argument of \(\propref{c14-pbs}\).
  \item \textbf{Tie direction.} Within a band, hosts are chosen by \emph{smallest
        index}; the sort of \(\mathit{ord}\) is \((\text{count},\text{city})\)
        ascending, and \(\mathcal{B}_{\mathit{idx}}\) indexes by city so the lift
        already returns the smallest. Sorting by city alone breaks the level.
  \item \textbf{\(\mathit{lowbit}\) and \(k\) are \code{long long}.} \(j\) and the
        prefix target exceed \(2^{31}\); the offset into the band still fits
        \(m\le 5\cdot10^5\), but \(S(L)\) and the running \(\mathit{acc}\) do not.
  \item \textbf{I/O.} \(n,q\le 5\cdot10^5\): \code{scanf}/fast I/O; sorting the
        queries by \(k\) (not reading order) is what makes the sweep monotone.
\end{itemize}
\end{hardeningbox}

\begin{recallbox}[Chapter 18 in ten seconds]
A Fenwick tree is an array hiding a binary tree on the index axis: \(\mathit{lowbit}(i)=i\&(-i)\)
defines slot \(i\)'s responsibility interval, a down-walk reads a prefix, an
up-walk writes a point (\(\defref{c14-interval}\), \(\thmref{c14-query}\)), and a
\(\bigO{n}\) build hands blocks forward (\(\propref{c14-linear}\)). Change the
\emph{stored quantity} and you get a family: differences for range-add/point-read
(\(\defref{c14-diff}\)), two BITs for range-add/range-sum (\(\thmref{c14-twobit}\)),
frequencies for order statistics lifted down the tree (\(\thmref{c14-lift}\)), the
Euler array for subtrees and paths (\(\propref{c14-eulerrange}\)). Change the
\emph{order of operations} and you get offline sweeps: sort one axis, BIT the
other (\(\propref{c14-dom}\)), wrap it in parallel binary search for a threshold
(\(\propref{c14-pbs}\)) -- a BIT only appends, so no rollback.
\selftest{
  \item Give the two walks and the loop-variable sign on each. \hfill\lv{3}
  \item Why is a min-BIT append-only, and what does that force on decrements? \hfill\lv{4}
  \item In the \(k\)-th lift, is the descent test \(<k\) or \(\le k\), and what do you return? \hfill\lv{4}
  \item State the two BIT axes the Summit uses and why one needs no rollback. \hfill\lv{5}
}
\end{recallbox}

\chaptersection{Tool Kit Inventory}

The machine assembled through Chapter~18. Earlier chapters' rows are carried
forward; this chapter's are appended. Pointers only.

\begin{center}
\small
\renewcommand{\arraystretch}{1.25}
\begin{tabular}{@{}p{0.21\linewidth} l c p{0.25\linewidth} p{0.17\linewidth}@{}}
\toprule
\textbf{Tool} & \textbf{Since} & \textbf{Lv} & \textbf{Trigger phrase} & \textbf{Where} \\
\midrule
\multicolumn{5}{l}{\itshape Carried forward} \\
Euler tour \(\tin/\tout\) flatten & 5 & \lv{3} & ``a subtree is a contiguous range'' & \chref{c24} (\(\thmref{c24-range}\)) \\
Merge-sort inversion count & 1 & \lv{4} & ``offline divide-and-conquer twin'' & \chref{c01} \\
Rollback-DSU PBS scheduling & 17 & \lv{5} & ``binary-search all queries in lockstep'' & \secref{c28-offline} \\
\midrule
\multicolumn{5}{l}{\itshape Chapter 18 -- new this chapter} \\
Plain BIT (query/update/build) & 18.1 & \lv{2} & ``live prefix sum, point edit'' & \coderef{c14-query}, \coderef{c14-build} \\
Difference BIT & 18.2 & \lv{3} & ``add to a range, read a point'' & \coderef{c14-diff-build} \\
Two-BIT range sum & 18.2 & \lv{4} & ``add to a range, sum a range'' & \coderef{c14-twobit} \\
Coordinate compression & 18.2 & \lv{3} & ``values sparse / up to 1e9'' & \coderef{c14-compress} \\
Frequency BIT + \(k\)-th lift & 18.3 & \lv{4} & ``select the \(k\)-th live item'' & \coderef{c14-kth} \\
Inversion / smaller-after sweep & 18.3 & \lv{4} & ``count inversions online'' & \coderef{c14-inv} \\
Euler-driven BIT (subtree/path) & 18.4 & \lv{4} & ``subtree add/sum on a tree'' & \coderef{c14-subtree}, \coderef{c14-path} \\
2D dominance / nested sweep & 18.5 & \lv{4} & ``sort one axis, BIT the other'' & \coderef{c14-dom}, \coderef{c14-nested} \\
PBS over a BIT & 18.5 & \lv{5} & ``least threshold, prefix predicate'' & \coderef{c14-pbs} \\
Count-space + index-space lift & Summit & \lv{5} & ``find a level, pick the offset-th'' & \coderef{c14-summit} \\
\bottomrule
\end{tabular}
\end{center}

\chaptersection{Related Topics}

\begin{itemize}[nosep, leftmargin=1.4em]
  \item \chref{c24} (Ch.~5, Euler Tour Technique) --- backward: owns the
        \(\tin/\tout\) flattening and the ``subtree \(=\) interval'' fact
        (\(\thmref{c24-range}\)); 18.4 reuses them by label and never re-derives
        them -- this chapter supplies the \emph{range engine} Ch.~5 promised, and
        its CSES~1137 practice is re-solved here with a BIT (\(\probref{c14-prop}\)).
  \item \chref{c28} (Ch.~17, DSU \& Sack) --- backward/forward: the parallel
        binary-search \emph{scheduling} is theirs (\(\secref{c28-offline}\)); 18.5
        discharges Ch.~17's forward pointer by swapping the rollback-DSU predicate
        for a BIT sweep (\(\propref{c14-pbs}\)). The Summit is the joint
        composition the handoff advertised.
  \item \chref{c01} (Ch.~1, complexity \& counting) --- backward: owns the
        merge-sort inversion count; 18.3 is the online frequency-BIT twin
        (\(\propref{c14-inversion}\)) -- same \(\bigO{n\log n}\), different mechanics.
  \item \chref{c06} (Ch.~6, Constructing Trees) --- back-cite: owns CF~1354D
        Multiset and PBDS order-statistics trees; 18.3 is the array-index rival
        (prose only, URLs not re-linked).
  \item \chref{c13} (Ch.~19, Segment Trees) --- forward: the BIT-vs-segtree
        boundary is drawn in 18.2 -- a BIT wins on prefix sums and simplicity, a
        segtree owns arbitrary-range min/non-invertible combines and lazy; no
        segtree code appears here.
  \item \chref{c15} (Ch.~20, Lazy Propagation) --- forward: range-add/range-sum
        with a \emph{third} operation (or non-commutative combine) belongs there,
        not to the two-BIT trick (\(\thmref{c14-twobit}\)'s remark).
\end{itemize}

\begin{lastword}[Last Word --- Chapter 18]
\emph{``One \(\mathit{lowbit}\), two walks, and an array starts answering prefix
questions in \(\bigO{\log n}\). Change what you store -- sums, differences,
frequencies, Euler slots -- and the same two walks become a whole family; change
the order you touch them -- sort one axis, sweep the other -- and offline
questions fall to the same structure.''}

\medskip
The plain BIT is honest and near-free (\(\coderef{c14-query}\)); the \(\bigO{n}\)
build makes it production-cheap (\(\propref{c14-linear}\)). Store differences and
a range add is two writes (\(\defref{c14-diff}\)); store a position-weighted pair
and range sums follow (\(\thmref{c14-twobit}\)). Index a frequency array and the
down-walk becomes a monotone search, so lifting the tree selects the \(k\)-th live
item (\(\thmref{c14-lift}\)) -- the reflex the Summit spends twice, once on a count
axis and once on a city axis (\(\probref{c14-summit}\)). Drive Ch.~5's flat array
and subtrees/paths are just ranges again (\(\propref{c14-eulerrange}\)); sort one
axis and sweep the other and dominance, nested ranges, and a threshold search are
one shared BIT with no undo (\(\propref{c14-dom}\), \(\propref{c14-pbs}\)).

\medskip
Carry three reflexes forward:
\begin{enumerate}[nosep, leftmargin=1.6em]
  \item \emph{Index \(=\) meaning.} A BIT slot can be an array position, a rank, an
        Euler entry, or a count value -- naming which is half the solution
        (\(\figref{c14-taxonomy}\)).
  \item \emph{Append-only is the BIT's contract.} Prefix sums and monotone sweeps
        fit; arbitrary decrements and non-invertible combines are the segment
        tree's job (\(\propref{c14-prefixmin}\), forward to \chref{c13}).
  \item \emph{Offline beats overkill.} If you may sort, a single sorted sweep plus
        a plain BIT replaces rollback stacks and persistent structures
        (\(\secref{c14-offline}\)).
\end{enumerate}

\medskip
And three habits, drilled until boring: keep the query and update walks'
loop-signs straight and 1-indexed; re-check the \(\tout\) inclusivity and the
\(r{+}1\) write against the smallest sample that could expose them; and before
reaching for a segment tree, ask whether the combine is a prefix sum -- if so this
chapter's eight lines already solve it.

\medskip
\textbf{Next:} \chref{c13} (Segment Trees) takes the range engine you now own and
generalises the combine -- arbitrary intervals, min/max with deletions, lazy
propagation -- paying a constant factor for the extra power. The BIT-vs-segtree
line drawn in \secref{c14-range} is exactly where that chapter picks up.
\end{lastword}
